\begin{theorem}[The ordinary direct images]\hypertarget{tag-02YN}{}\label{night:P5.2}\label{stmt:P5.2}
\StatementRecord{night:P5.2}{proved}{02YN}{\cref{inv:D1.1,inv:D1.2,push:src:E.3.1,night:P5.1}}
For the sheaves of \cref{I:ordinary-tangent-sheaves}, the coarse
direct images are
\[
\begin{array}{c|ccc}
 & f_* & R^1f_* &R^2f_*\\ \hline
 \mathcal O_X&\mathcal O&\mathcal O\oplus\mathcal O(-1)&\mathcal O(-1)\\
 L&\mathcal O(-1)&\mathcal O(-1)^{\oplus2}&\mathcal O(-1)\\
 F&\mathcal O&\mathcal O^{\oplus2}&\mathcal O\\
 \mathcal V&\mathcal O\oplus\mathcal O(-1)&\mathcal O\oplus\mathcal O(-1)^{\oplus3}
       &\mathcal O(-1)^{\oplus2}\\
 J&\mathcal O(-1)&\mathcal O(-1)^{\oplus2}&\mathcal O .
\end{array}
\]
\end{theorem}
\begin{proof}
\begin{enumerate}[label=\arabic*.,leftmargin=*]
\item The root splitting and structure-sheaf direct images are given in
\cref{inv:D1.2,push:src:E.3.1}; in the standard data these are also computed in
\cite[Proposition~6.6, pp.~77--78; Lemma~6.7, p.~78;
Proposition~6.10, pp.~79--80]{DuEtAl2026v1}.
The determinant line in these conventions is
\begin{equation}\hypertarget{tag-01OG}{}\label{eq:P5L}
 \det E=\widehat L
 =\rho^*\mathcal O_B(-1)\otimes
   \mathcal O_{\mathcal B}(2p_1/3+p_2/4)
 =\mathcal O_{\mathcal B}(-1).
\end{equation}
The equality is the tensor product of the divisor-section maps
in \cref{I:ordinary-tangent-sheaves}; its degree is $-12+8+3$.
For integers $d,a,b$, the invariant root-line formula is
\begin{equation}\hypertarget{tag-01OI}{}\label{eq:P5floor}
 \rho_*\bigl(\rho^*\mathcal O_B(d)\otimes
 \mathcal O_{\mathcal B}(ap_1/3+bp_2/4)\bigr)
 =\mathcal O_B\bigl(d+\lfloor a/3\rfloor+\lfloor b/4\rfloor\bigr).
\end{equation}
Indeed writing $a=3q+r$, $b=4q'+r'$ with
$0\leq r<3$, $0\leq r'<4$ separates the coarse divisor
$(d+q+q')$ from the two root characters. This is the same
invariant-section map as \cref{inv:D1.1}, and applies to all
integer $a,b$.

\item For a root line bundle $A$, projection and exact finite
invariants give
\[
 R^qf_*\bar f^*A=\rho_*\bigl(A\otimes
              \bigwedge^q(\mathcal O\oplus\widehat L)\bigr).
\]
For $A=\mathcal O,\widehat L,\widehat F$, and for
$E=\mathcal O\oplus\widehat L$, the formula
\eqref{eq:P5floor} gives the first four rows. The root-to-coarse
map, including the full scheme fibres, is \cref{I:root-coarse-comparison}.

\item For the last row one must compute restriction to $D_{\mathrm{dbl}}$.  The conductor
square gives
\[
0\to\mathcal O_W\to\nu_*\mathcal O_S\oplus\mathcal O_{D_{\mathrm{dbl}}}
\to\nu_*\mathcal O_{\partial S}\to0,
\]
where $\partial S$ is the hexagon.  Here $h^1(\mathcal O_{\partial S})=1$
and $h^1(\mathcal O_{D_{\mathrm{dbl}}})=2$.  The equality conditions at the two triple
points give incidence matrix
\[
 \begin{pmatrix}1&-1&0\\0&1&-1\\1&-1&0\\0&1&-1\end{pmatrix}.
\]
Orient the three theta-graph edges from $P$ to $Q$.  The hexagon
traversal maps to $e_1-e_2+e_3-e_1+e_2-e_3=0$.  The constant-to-coherent
branch comparison is an isomorphism in degree one, so the induced map
$H^1(D_{\mathrm{dbl}},\mathcal O_{D_{\mathrm{dbl}}})\to H^1(\partial S,\mathcal O_{\partial S})$ is
zero.  The conductor sequence then makes
$H^1(W,\mathcal O_W)\to H^1(D_{\mathrm{dbl}},\mathcal O_{D_{\mathrm{dbl}}})$ an isomorphism.

\item The canonical section of $F$ is nonzero near $W$, so $F|_{D_{\mathrm{dbl}}}=\mathcal O_{D_{\mathrm{dbl}}}$.
In $0\to J\to F\to\mathcal O_{D_{\mathrm{dbl}}}\to0$, the maps of direct images
$\mathcal O\to\mathbb C_{p_0}$ and
$\mathcal O^{\oplus2}\to\mathbb C_{p_0}^{\oplus2}$ are evaluation and
evaluation followed by this isomorphism.  They are surjective, with
kernels $\mathcal O(-p_0)$ in the indicated ranks.  Finally
$H^2(D_{\mathrm{dbl}},\mathcal O_{D_{\mathrm{dbl}}})=0$ and $R^2f_*F=\mathcal O$.  This gives the last row.
\end{enumerate}
\end{proof}